\tocsection{\textbf{11/10/26 Lecture 20:} $L^p$-Convergence and Random Series}

\begin{lemma}[\textbf{Proposition 20.1:} $L^p$-Convergence]
    Fix $p > 0$ and let $\xi, \xi_1, \xi_2, \ldots \in L^p$ with $\xi_n \xrightarrow{P} \xi$. Then, these conditions are equivalent:
    \begin{enumerate}[label = (\roman*)]
        \item $\xi_n \to \xi$ in $L^p$
        \item $\norm{\xi_n}_p \to \norm{\xi}_p$
        \item The variables $\abs{\xi_n}^p$, $n \in \N$ are uniformly integrable
    \end{enumerate}
    Conversely, (i) implies $\xi_n \xrightarrow{P} \xi$.
\end{lemma}
\begin{proof}{Proof:}
    $(\Leftarrow)$ Assume $\xi_n \to \xi$ in $L^p$. Then $\norm{\xi_n}_p \to \norm{\xi}_p$ by a previous result, and we see that, for any $\epsilon > 0$,
    \[ P\{\abs{\xi_n - \xi} > \epsilon\} = P\{\abs{\xi_n - \xi}^p > \epsilon^p\} \leq \epsilon^{-p}\norm{\xi_n - \xi}_p^p \to 0 \]
    by Chebychev, so $\xi_n \xrightarrow{P} \xi$.\\

    $(\Rightarrow)$ Now, assume $\xi_n \xrightarrow{P} \xi$. Then, $\abs{\xi_n}^p \xrightarrow{d} \abs{\xi}^p$ by the continuous mapping lemma and as convergence in probability implies convergence in distribution, so by Lemma 19.2, (ii) and (iii) are equivalent.\\

    
    Now, assume (ii) and suppose (i) fails. Then, there exists a subsequence $N' \subseteq \N$ such that $\inf_{n \in N'} \norm{\xi_n - \xi}_p > 0$. Now, as $\xi_n \xrightarrow{P} \xi$, we can find a further subsequence $N'' \subseteq N'$ such that $\xi_n \to \xi$ a.s.\ along $N''$, which implies $\norm{\xi_n - \xi}_p \to 0$ along $N''$, which is a contradiction.\\

    Thus, (ii) $\implies$ (i), and as (i) $\implies$ (ii), we are done.

    $\qed$
\end{proof}

\begin{lemma}[\textbf{Proposition 20.2:} Sums with Positive Terms]
    Let $\xi_1, \xi_2, \ldots$ be independent $\R_{+}$-valued random variables. Then, $\sum_n \xi_n < \infty$ a.s.\ iff $\sum_n \mathbb{E}[\xi_n \wedge 1] < \infty$.
\end{lemma}
\begin{proof}{Proof:}
    Assuming the condition, we see $\mathbb{E}\sum_n(\xi_n \wedge 1) < \infty$ by Fubini's theorem, so $\sum_n(\xi_n \wedge 1) < \infty$ a.s.\\

    In particular, $\sum_n 1\{\xi_n > 1\} < \infty$ a.s., so $\sum_n(\xi_n \wedge 1)$ and $\sum_n \xi_n$ differ at finitely many values. Thus, $\sum_n \xi_n < \infty$ a.s.\\

    Conversely, assume $\sum_n \xi_n < \infty$ a.s. Then, also $\sum_n(\xi_n \wedge 1) < \infty$ a.s., so we can assume WLOG $\xi_n \leq 1$ for all $n$.\\

    Using that $1 - x \leq e^{-x} \leq 1 - ax$ for $x \in [0, 1]$, where $a = 1 - e^{-1}$,
    \begin{align*}
        0 < \mathbb{E}\left[e^{-\sum_n \xi_n}\right] &= \prod_n \mathbb{E}[e^{-\xi_n}]\\
        &\leq \prod_n(1 - a\mathbb{E}[\xi_n])\\
        &\leq \prod_n e^{-a\mathbb{E}[\xi_n]}\\
        &= e^{-a\sum_n \mathbb{E}[\xi_n]}
    \end{align*}
    (the first inequality holds since $\sum_n \xi_n < \infty$ a.s.), and so $\sum_n \mathbb{E}\xi_n < \infty$, as otherwise $0 = e^{-a\sum_n \mathbb{E}\xi_n}$.

    $\qed$
\end{proof}

\textbf{Note:} $\sum_n \xi_n$ converges $\iff \lim_{N \to \infty} \sup_{m} \abs{\xi_{N + 1} + \cdots + \xi_{N + m}} = 0$.

\begin{lemma}[\textbf{Lemma 20.3:} Maximum Inequality (Kolmogorov)]
    Let $\xi_1, \xi_2, \ldots$ be independent random variables with mean zero, and put $S_n = \xi_1 + \cdots + \xi_n$. Then,
    \[ P\left\{\sup_n \abs{S_n} > r\right\} \leq r^{-2}\sum_n \mathbb{E}\xi_n^2, \quad r > 0 \]
\end{lemma}
\begin{proof}{Proof:}
    Assume $\sum_n \mathbb{E}\xi_n^2 < \infty$ (otherwise obvious). Letting $\tau = \inf\{n \col \abs{S_n} > r\}$, and noting that
    \[ S_k1\{\tau = k\} \perp\!\!\!\perp (S_n - S_k) \quad \tt{for } k \leq n \]
    since $S_k1\{\tau = k\}$ is measurable with respect to $\sigma(\xi_1, \ldots, \xi_k)$ and $S_n - S_k$ is measurable with respect to $\sigma(\xi_{k + 1}, \ldots, \xi_n)$. Then, we see
    \begin{align*}
        \sum_{k \leq n} \mathbb{E}\xi_k^2 &= \mathbb{E}S_n^2 \quad (\mathbb{E}\xi_n = 0)\\
        &\geq \sum_{k \leq n} \mathbb{E}[S_n^2 ; \tau = k]\\
        &\geq \sum_{k \leq n}\left\{\mathbb{E}[S_k^2 ; \tau = k] + 2\mathbb{E}[S_k(S_n - S_k) ; \tau = k]\right\}\\
        &= \sum_{k \leq n} \mathbb{E}[S_k^2 ; \tau = k]\\
        &\geq \sum_{k \leq n} r^2P\{\tau = k\}\\
        &= r^2P\{\tau \leq n\}
    \end{align*}
    where
    \begin{itemize}
        \item the second inequality uses $S_n^2 = S_k^2 + 2S_k(S_n - S_k) + (S_n - S_k)^2$,
        \item the equality uses $S_k1\{\tau = k\} \perp\!\!\!\perp S_n - S_k$ and $\mathbb{E}[S_n - S_k] = 0$,
        \item the last inequality uses $S_k^2 1\{\tau = k\} \geq r^2 1\{\tau = k\}$.
    \end{itemize}
    Taking $n \to \infty$,
    \[ \sum_k \mathbb{E}\xi_k^2 \geq r^2P\{\tau < \infty\} = r^2P\left\{\sup_k \abs{S_k} > r\right\} \]

    $\qed$
\end{proof}

\textbf{Remark:} This is quite remarkable since
\[ P\{\abs{S_n} > r\} \leq \frac{\mathbb{E}[S_n^2]}{r^2} = \frac{\sum_{k \leq n} \mathbb{E}\xi_k^2}{r^2} \quad (\tt{mean } 0) \]
so Chebyshev extends!

\begin{lemma}[\textbf{Lemma 20.4:} Variance Criterion for Series]
    Let $\xi_1, \xi_2, \ldots$ be independent random variables with mean $0$ and $\sum_n \mathbb{E}\xi_n^2 < \infty$. Then, $\sum_n \xi_n$ converges a.s.
\end{lemma}
\begin{proof}{Proof:}
    Write $S_n = \xi_1 + \cdots + \xi_n$. By the lemma, for any $\epsilon > 0$,
    \[ P\left\{\sup_{k \geq n} \abs{S_n - S_k} > \epsilon\right\} \leq \epsilon^{-2}\sum_{k \geq n} \mathbb{E}\xi_k^2 \]
    Thus, $\sup_{k \geq n} \abs{S_n - S_k} \xrightarrow{P} 0$ as $n \to \infty$. Then, $\sup_{k \geq n} \abs{S_n - S_k} \to 0$ a.s.\ along a subsequence.\\

    By the triangle inequality,
    \[ \sup_{k, \ell \geq n} \abs{S_k - S_{\ell}} \leq \sup_{k \geq n} \abs{S_n - S_k} + \sup_{\ell \geq n} \abs{S_n - S_{\ell}} \]
    so the LHS converges to $0$ along the subsequence. Moreover, the sup is non-increasing in $n$, so a.s.\ convergence extends to the entire sequence. Thus, $(S_n)$ is a.s.\ Cauchy.\\

    Thus, $S_n$ converges a.s.\ by a previous lemma.

    $\qed$
\end{proof}

\begin{theorem}[\textbf{Theorem 20.5:} Positive and Symmetric Terms]
    Let $\xi_1, \xi_2, \ldots$ be independent symmetric random variables. Then, these conditions are equivalent:
    \begin{enumerate}[label = (\roman*)]
        \item $\sum_n \xi_n$ converges a.s.
        \item $\sum_n \xi_n^2 < \infty$ a.s.
        \item $\sum_n \mathbb{E}[\xi_n^2 \wedge 1] < \infty$
    \end{enumerate}
    If these conditions fail, then $\abs{\sum_{k \leq n} \xi_k} \xrightarrow{P} \infty$; that is, for all $r > 0$,
    \[ P\left\{\abs{\sum_{k \leq n} \xi_k} < r\right\} \to 0 \]
\end{theorem}
\begin{proof}{Proof:}
    We have (ii) $\iff$ (iii) from earlier.\\

    
  
    Now, assume (iii). Then, $\sum_n \xi_n1\{\abs{\xi_n} \leq 1\}$ converges a.s.\ by Lemma 20.4 as $\mathbb{E}[\xi_n1\{\abs{\xi_n} \leq 1\}] = 0$ and
    \[ \sum_n \mathbb{E}\left[(\xi_n1\{\abs{\xi_n} \leq 1\})^2\right] \leq \sum_n \mathbb{E}[\xi_n^2 \wedge 1] < \infty \]
    Then, by Fubini's theorem applied to (iii), $\sum_n 1\{\abs{\xi_n} > 1\} < \infty$ a.s., so $\sum_n \xi_n1\{\abs{\xi_n} \leq 1\}$ and $\sum_n \xi_n$ differ by finitely many terms, so $\sum_n \xi_n$ converges a.s.\\

    Now, assume (i) and suppose (ii) fails. Then, $\sum_n \xi_n^2 = \infty$ a.s.\ by Kolmogorov's 0-1 law. Using an argument outside the scope of this course, we see $\abs{S_n} \xrightarrow{P} \infty$. Thus, $\abs{S_n} \to \infty$ a.s.\ along a subsequence, which contradicts (i).

    $\qed$
\end{proof}
